\NormalFont\ShortTitle{Coloseni}

{\MT Scrisoarea lui Pavel către Coloseni


\par }\ChapOne{1}{\PP \VerseOne{1}Pavel, apostol al lui Isus Hristos, prin voia lui Dumnezeu, și Timotei, fratele nostru,

\VS{2}către sfinții și frații credincioși în Hristos din Colose: Harul și pacea să vă fie vouă, de la Dumnezeu Tatăl nostru și de la Domnul Isus Hristos.


\par }{\PP \VS{3}Mulțumim lui Dumnezeu, Tatăl Domnului nostru Isus Hristos, rugându-ne totdeauna pentru voi,

\VS{4}după ce am auzit de credința voastră în Hristos Isus și de dragostea pe care o aveți față de toți sfinții,

\VS{5}din pricina nădejdii care vă este păstrată în ceruri, despre care ați auzit mai înainte în cuvântul adevărului bunei vestiri

\VS{6}care a ajuns la voi, așa cum este în toată lumea și care dă roade și crește, așa cum este și în voi, din ziua în care ați auzit și ați cunoscut harul lui Dumnezeu în adevăr,

\VS{7}așa cum ați aflat de la Epafras, iubitul nostru tovarăș de slujbă, care este un slujitor credincios al lui Hristos în numele vostru,

\VS{8}care ne-a și vestit dragostea voastră în Duhul Sfânt.


\par }{\PP \VS{9}De aceea și noi, din ziua când am auzit aceasta, nu încetăm să ne rugăm și să facem cereri pentru voi, ca să fiți plini de cunoștința voii Lui în toată înțelepciunea și priceperea spirituală,

\VS{10}ca să umblați în chip vrednic de Domnul, ca să-I plăce în toate privințele, aducând roade în orice lucrare bună și sporind în cunoștința lui Dumnezeu,

\VS{11}întăriți cu toată puterea, după puterea slavei sale, pentru toată stăruința și perseverența cu bucurie,

\VS{12}mulțumind Tatălui, care ne-a făcut apți să fim părtași la moștenirea sfinților în lumină,

\VS{13}care ne-a izbăvit din puterea întunericului și ne-a transferat în Împărăția Fiului iubirii sale,

\VS{14}în care avem răscumpărarea noastră, iertarea păcatelor noastre.


\par }{\PP \VS{15}El este chipul lui Dumnezeu cel nevăzut, întâiul născut al întregii creații.

\VS{16}Căci prin el au fost create toate lucrurile în ceruri și pe pământ, cele văzute și cele nevăzute, fie tronuri, fie stăpâniri, fie principate, fie puteri. Toate lucrurile au fost create prin el și pentru el.

\VS{17}El este înaintea tuturor lucrurilor și în El toate lucrurile sunt ținute împreună.

\VS{18}El este capul trupului, al adunării, El este începutul, Cel întâi născut din morți, pentru ca în toate lucrurile să aibă întâietatea.

\VS{19}Căci toată plinătatea a binevoit să locuiască în el,

\VS{20}și prin el să împace toate lucrurile cu sine prin el, fie cele de pe pământ, fie cele din ceruri, făcând pace prin sângele crucii sale.


\par }{\PP \VS{21}Pe voi, care în trecut erați înstrăinați și vrăjmași în mintea voastră prin faptele voastre rele,

\VS{22}dar acum v-a împăcat în trupul trupului Său, prin moarte, ca să vă înfățișeze înaintea Lui sfinți, fără cusur și fără prihană,

\VS{23}dacă rămâneți în credință, întemeiați și neclintiți, fără să vă abateți de la nădejdea Bunei Vestiri, pe care ați auzit-o, și care este vestită în toată făptura de sub ceruri, și de care eu, Pavel, am fost făcut slujitor.


\par }{\PP \VS{24}Și eu mă bucur de suferințele mele pentru voi, și împlinesc ceea ce-mi lipsește din suferințele lui Hristos în trupul meu, pentru trupul Lui, care este Adunarea,

\VS{25}a cărei slujbă am fost făcut slujitor, potrivit cu slujba lui Dumnezeu, care mi-a fost dată față de voi, ca să împlinesc cuvântul lui Dumnezeu,

\VS{26}taina care a fost ascunsă de veacuri și de generații. Dar acum a fost descoperit sfinților Săi,

\VS{27}cărora Dumnezeu a binevoit să le facă cunoscut care sunt bogățiile slavei acestui mister printre neamuri, care este Hristos în voi, nădejdea slavei.

\VS{28}Noi îl vestim, sfătuind pe fiecare om și învățând pe fiecare în toată înțelepciunea, ca să prezentăm pe fiecare om desăvârșit în Hristos Isus;

\VS{29}pentru care și eu mă străduiesc, străduindu-mă după lucrarea Lui, care lucrează în mine cu putere.



\par }\Chap{2}{\PP \VerseOne{1}Căci vreau să știți cât de mult mă lupt pentru voi și pentru cei din Laodiceea, și pentru toți cei care nu Mi-au văzut fața în trup,

\VS{2}ca să fie mângâiate inimile lor, fiind uniți în dragoste și dobândind toate bogățiile unei înțelegeri depline, ca să cunoască taina lui Dumnezeu, atât a Tatălui, cât și a lui Hristos,

\VS{3}în care sunt ascunse toate comorile înțelepciunii și ale cunoștinței.

\VS{4}Spun aceasta pentru ca nimeni să nu vă amăgească cu persuasiunea discursului.

\VS{5}Căci, chiar dacă sunt absent în trup, sunt cu voi în duh, bucurându-mă și văzând ordinea voastră și statornicia credinței voastre în Hristos.


\par }{\PP \VS{6}Deci, după cum ați primit pe Hristos Isus, Domnul, umblați în El,

\VS{7}înrădăcinați și zidiți în El și întăriți în credință, așa cum ați fost învățați, sporind în ea cu mulțumiri.


\par }{\PP \VS{8}Fiți atenți să nu vă lăsați jefuiți de cineva prin filozofia lui și prin înșelăciune deșartă, după tradiția oamenilor, după duhurile elementare ale lumii, și nu după Hristos.

\VS{9}Căci în el locuiește trupește toată plinătatea dumnezeirii,

\VS{10}și în el sunteți desăvârșiți, el care este capul tuturor principatelor și puterilor.

\VS{11}În El ați fost și voi tăiați împrejur cu o circumcizie care nu se face cu mâinile, prin lepădarea trupului de păcatele cărnii, în circumcizia lui Hristos,

\VS{12}fiind îngropați împreună cu El în botez, în care ați și înviat împreună cu El, prin credința în lucrarea lui Dumnezeu, care L-a înviat din morți.

\VS{13}Voi erați morți prin greșelile voastre și prin necircumcizia cărnii voastre. El v-a făcut vii împreună cu el, după ce ne-a iertat toate greșelile noastre,

\VS{14}ștergând scrisul de mână din rânduieli care era împotriva noastră. El l-a scos din calea noastră, bătându-l în cuie pe cruce.

\VS{15}După ce a dezbrăcat principatele și puterile, le-a dat în vileag, triumfând asupra lor în ea.


\par }{\PP \VS{16}Nimeni, deci, să nu vă judece în privința mâncării sau a băuturii, sau în privința sărbătorilor, a lunii noi sau a Sabatului,

\VS{17}care sunt o umbră a lucrurilor viitoare; dar trupul este al lui Hristos.

\VS{18}Nimeni să nu vă răpească premiul vostru prin înjosirea de sine și închinarea la îngeri, stăruind în lucrurile pe care nu le-a văzut, îngâmfat în zadar de mintea lui trupească,

\VS{19}și neținându-se ferm de Cap, de la care tot trupul, fiind alimentat și unit prin încheieturi și ligamente, crește cu creșterea lui Dumnezeu.


\par }{\PP \VS{20}Dacă ați murit împreună cu Hristos de duhurile elementare ale lumii, de ce, ca și cum ați trăi în lume, vă supuneți la rânduieli,

\VS{21}“Nu mânuiți, nu gustați, nu atingeți”

\VS{22}(care toate pier cu folosul), după poruncile și învățăturile oamenilor?

\VS{23}Aceste lucruri, într-adevăr, par a fi înțelepciune în ceea ce privește cultul autoimpus, umilința și severitatea față de trup, dar nu au nicio valoare împotriva desfătării cărnii.



\par }\Chap{3}{\PP \VerseOne{1}Deci, dacă ați fost înviați împreună cu Hristos, căutați lucrurile de sus, unde este Hristos, așezat la dreapta lui Dumnezeu.

\VS{2}Puneți-vă mintea la lucrurile de sus, nu la cele de pe pământ.

\VS{3}Căci ați murit, iar viața voastră este ascunsă cu Hristos în Dumnezeu.

\VS{4}Când Hristos, viața noastră, se va arăta, atunci și voi veți fi arătați împreună cu el în glorie.


\par }{\PP \VS{5}Omorâți, așadar, membrele voastre care sunt pe pământ: imoralitatea sexuală, necurăția, pasiunea depravată, dorința rea și lăcomia, care este idolatrie.

\VS{6}Din pricina acestor lucruri, mânia lui Dumnezeu vine peste copiii neascultării.

\VS{7}Și voi ați umblat odinioară în ele, când trăiați în ele,

\VS{8}dar acum trebuie să le îndepărtați pe toate: mânia, furia, răutatea, calomnia și vorbirea rușinoasă din gura voastră.

\VS{9}Nu vă mințiți unii pe alții, văzând că v-ați dezbrăcat de omul cel vechi cu faptele lui,

\VS{10}și v-ați îmbrăcat în omul cel nou, care se înnoiește în cunoaștere, după chipul Creatorului său,

\VS{11}unde nu poate fi vorba de greci și iudei, de circumcizie și necircumcizie, de barbari, de sciți, de robi sau de oameni liberi, ci Hristos este totul și în toți.


\par }{\PP \VS{12}Îmbrăcați-vă deci, ca niște aleși ai lui Dumnezeu, sfinți și iubiți, o inimă plină de compasiune, de bunătate, de smerenie, de umilință și de stăruință;

\VS{13}purtați-vă unii pe alții și iertați-vă unii pe alții, dacă are cineva vreo plângere împotriva cuiva; după cum v-a iertat Hristos, așa să faceți și voi.


\par }{\PP \VS{14}Mai presus de toate aceste lucruri, umblați în dragoste, care este legătura desăvârșirii.

\VS{15}Și să domnească în inimile voastre pacea lui Dumnezeu, la care ați și fost chemați într-un singur trup, și fiți recunoscători.

\VS{16}Cuvântul lui Hristos să locuiască bogat în voi; cu toată înțelepciunea, învățându-vă și îndemnându-vă unii pe alții cu psalmi, imnuri și cântări duhovnicești, cântând cu har în inima voastră Domnului.


\par }{\PP \VS{17}Orice faceți, cu cuvântul sau cu fapta, faceți totul în Numele Domnului Isus, mulțumind prin El lui Dumnezeu Tatăl.


\par }{\PP \VS{18}Neveste, supuneți-vă bărbaților voștri, cum se cuvine în Domnul.


\par }{\PP \VS{19}Bărbați, iubiți-vă nevestele și nu vă înverșunați împotriva lor.


\par }{\PP \VS{20}Copii, ascultați de părinții voștri în toate lucrurile, căci aceasta place Domnului.


\par }{\PP \VS{21}Părinți, nu vă provocați copiii, ca să nu se descurajeze.


\par }{\PP \VS{22}Robilor, ascultați în toate lucrurile de cei ce vă sunt stăpâni după trup, nu numai când vă privesc, ca niște plăceri, ci cu inimă curată, temându-vă de Dumnezeu.

\VS{23}Și orice faceți, lucrați din toată inima, ca pentru Domnul și nu pentru oameni,

\VS{24}știind că de la Domnul veți primi răsplata moștenirii, căci voi slujiți Domnului Hristos.

\VS{25}Dar cine greșește, va primi înapoi pentru greșeala pe care a făcut-o, și nu este nici o părtinire.



\par }\Chap{4}{\PP \VerseOne{1}Stăpâni, dați robilor voștri ceea ce este drept și egal, știind că și voi aveți un Stăpân în ceruri.


\par }{\PP \VS{2}Continuați cu stăruință în rugăciune, veghind în ea cu mulțumire,

\VS{3}rugându-vă împreună și pentru noi, ca Dumnezeu să ne deschidă o ușă pentru cuvânt, ca să spunem taina lui Hristos, pentru care și eu sunt în legături,

\VS{4}ca să o pot descoperi așa cum trebuie să vorbesc.


\par }{\PP \VS{5}Umblați cu înțelepciune față de cei de afară, răscumpărând timpul.

\VS{6}Vorbirea voastră să fie întotdeauna cu har, condimentată cu sare, ca să știți cum trebuie să răspundeți fiecăruia.


\par }{\PP \VS{7}Toate treburile mele vă vor fi aduse la cunoștință prin Ticus, fratele preaiubit, slujitorul credincios și tovarășul de robie în Domnul.

\VS{8}Tocmai de aceea vi-l trimit la voi, ca să vă cunoască situația și să vă mângâie inimile,

\VS{9}împreună cu Onesimus, fratele credincios și iubit, care este unul dintre voi. Ei vă vor face cunoscut tot ceea ce se întâmplă aici.


\par }{\PP \VS{10}Vă salută Aristarh, tovarășul meu de închisoare, și Marcu, vărul lui Barnaba, despre care ați primit porunca: “Dacă vine la voi, primiți-l”,

\VS{11}și Isus, numit Iustus. Aceștia sunt singurii mei tovarăși de muncă pentru Împărăția lui Dumnezeu care sunt din circumcizie, oameni care mi-au fost o mângâiere.


\par }{\PP \VS{12}Epafras, care este unul dintre voi, slujitor al lui Hristos, vă salută și se străduiește mereu în rugăciunile sale pentru voi, ca să fiți desăvârșiți și să fiți desăvârșiți în toată voia lui Dumnezeu.

\VS{13}Căci mărturisesc despre el că are un mare zel pentru voi, pentru cei din Laodiceea și pentru cei din Ierapole.

\VS{14}Luca, medicul iubit, și Demas vă salută.

\VS{15}Salutați-i pe frații care sunt în Laodiceea, împreună cu Nimfa și cu adunarea care este în casa lui.

\VS{16}După ce se va citi această scrisoare între voi, faceți să fie citită și în adunarea laodiceenilor și să citiți și voi scrisoarea de la Laodiceea.

\VS{17}Spune-i lui Arhipus: “Ai grijă de slujba pe care ai primit-o în Domnul, ca să o împlinești.”


\par }{\PP \VS{18}Eu, Pavel, scriu cu mâna mea această salutare. Amintiți-vă de lanțurile mele. Harul să fie cu voi. Amin.


\par }